\scanpage{50}\mspage{10}
\[\unsure{\psi_{\vec a}|u_{\vec a}|\vec\psi_{\vec a}}=\unsure{u_{\vec b'}|v_{\vec b'}|\psi_{\vec b'}}u_{\vec b}\tag*{(34)}\]
où $\vec b$ est $b$ orienté de $s$ vers $t$, $\vec b'$ est $b'$ orienté
de $t$ vers $s'$, et $u_{\vec b'}$ et $u_{\vec b}$ sont définis
comme composés des $u_{\vec c}$ correspondants
aux arcs \inserted{de $K$} composants des chemins $\overrightarrow{(s,t)}$ et $\overrightarrow{(t,s')}$.
Il faudrait, dans cette optique, poser un
(A5) — savoir que ce $u_{\unsure{\vec a|\text{--}}}$ ne dépend pas des
choix arbitraires qui ont été faits. Mais
le cas où $K$ comporte une ps droite qui
est un bigône \deleted{\unsure{à mon avis il est plus}} \inserted{devrait plutôt être}
traité à part, avec des moyens plus
directement adaptés à cette situation. Je
n'ai donc pas vérifié que la proposition
% slanted
\margin{mais si, par l'argument en marge de la page précédente}
énoncée soit correcte, sans exclure le
cas où $K$ est formé de $n$ pseudo-droites
concourantes, coupées par une seule
autre ps. droite.

J'ai maintenant envie de montrer
que dans le cas d'un $K$\ldots \unsure{provient|provenant} d'un $\unsure{\mathcal X|X}$,
son orientation peut se décrire canoniquement
en termes de $(K,(D_i))$ seulement.
% archive sheet number: 49
